\section{Read counts before large position lists}
\label{sec:index-cold-loading}
\label{sec:index-format6}

Suppose the query contains hashes 12, 35, and 78. Reading every target posting
would waste work on keys that cannot produce an anchor for this query. It can
also waste work on a highly repeated key that the occurrence filter will
reject. The selective MPS format, called Format~6, separates the information
needed to make these two decisions.

The reader performs four steps:
\begin{enumerate}
  \item Read and check the complete directory, then find hash ranges
    intersecting the distinct query hashes.
  \item Read the selected controls to find each key's exact global
    occurrence count.
  \item Apply the occurrence threshold before requesting positions.
  \item Read accepted posting ranges and form the lookup needed by the
    alignment batch.
\end{enumerate}
``Global'' here means over the indexed target sequence. Looking at one loaded
piece alone could underestimate repetition and wrongly admit a frequent key.
The stored compact counts therefore have additional exact-count information
when their small field cannot hold the value.

\begin{figure}[htbp]
\centering
\begin{tikzpicture}[x=1cm,y=1cm,font=\footnotesize]
  \node[anchor=west,font=\bfseries] at (0,6.25) {Query hashes $\{12,35,78\}$; occurrence limit 5};
  \node[anchor=east] at (1.2,5.35) {directory};
  \foreach \x/\range/\col in {2.6/{0--19}/indexBlue,5.2/{20--39}/indexBlue,7.8/{40--59}/black,10.4/{60--79}/indexBlue}
    \node[index cell,fill=\col!10,draw=\col,minimum width=2.3cm] at (\x,5.35) {$\range$};
  \node[anchor=east] at (1.2,3.95) {controls};
  \node[index green,minimum width=2.3cm] (c12) at (2.6,3.95) {$12:\ count=2$};
  \node[index orange,minimum width=2.3cm] (c35) at (5.2,3.95) {$35:\ count=90$};
  \node[index cell,dashed,minimum width=2.3cm] at (7.8,3.95) {not read};
  \node[index green,minimum width=2.3cm] (c78) at (10.4,3.95) {$78:\ count=3$};
  \foreach \x in {2.6,5.2,10.4}
    \draw[index arrow] (\x,5.04)--(\x,4.28);
  \node[anchor=east] at (1.2,2.0) {postings};
  \draw[fill=indexGreen!16,draw=indexGreen] (1.45,1.65) rectangle (3.1,2.35);
  \node at (2.275,2.0) {$17,42$};
  \draw[fill=black!7,draw=black!45] (3.25,1.65) rectangle (8.9,2.35);
  \node at (6.075,2.0) {90 positions: not read};
  \draw[fill=indexGreen!16,draw=indexGreen] (9.05,1.65) rectangle (11.55,2.35);
  \node at (10.3,2.0) {$18,54,96$};
  \draw[index arrow,indexGreen] (c12.south)--(2.275,2.39);
  \draw[index arrow,indexGreen] (c78.south)--(10.3,2.39);
  \draw[thick,indexOrange] (5.2,3.61)--(5.2,2.92);
  \draw[thick,indexOrange] (5.02,2.92)--(5.38,2.92);
  \node[anchor=west,text=indexOrange] at (5.52,3.03) {$90>5$: stop before positions};
  \node[anchor=west,align=left] at (1.45,.6)
    {The index retains all occurrences.\\This query reads only the accepted lists.};
\end{tikzpicture}
\caption{Two decisions save different reads. All directory entries are
read; blue ranges select the controls. Counts reject a frequent hash before
its long list is loaded. Positions and thresholds are illustrative.}
\label{fig:index-selective-load}
\label{fig:format6-streams}
\end{figure}

Adjacent accepted ranges can be read together through a bounded I/O buffer.
The PostgreSQL representation stores payloads in heap rows, so a requested
range may fetch a complete row containing some additional bytes. Selective
logical ranges therefore reduce unnecessary reads without making the physical
read granularity equal to one posting.

The selective MPS reader prepares a lookup for the complete stored target
identified by its locator. This reader does not support transformed targets
or arbitrary SQL slices. Other index readers can clip and rebase positions
for supported intervals. The conditions for using each reader remain
separate from the lazy sequence interface itself.
